This proves c) and d) in the case of artinian $S$ (N.B. We shall prove below that in this case $\mathscr T$ is in fact affine over $S$).

Suppose now $S$ arbitrary. To prove c) and d), which are local assertions on $S$ for the Zariski topology, one can suppose that there exists a prime number $\ell$ prime to the residue characteristics of $S$, that $S$ is affine, and that the reductive rank of the fibers of $G$ (which is obviously locally constant, thanks to the hypothesis of local existence for fpqc of a maximal torus) is constant, say $r$. Let $S'\to S$ be a faithfully flat quasi-compact morphism, $S'$ affine, such that $G_{S'}$ admits a maximal torus $T'$. Let $C'$ be its centralizer; I say that there exists a suitable power $n$ of $\ell$ such that one also has $C=\Centr_G({}_nT)$. Indeed, to see this, one is reduced at once to the case of noetherian $S'$, where this was seen in XI~6.2. Fix $n$ thus; put
\[
M=(\ZZ/n\ZZ)^r,\qquad P=\uHom_{S\text{-gr}}(M_S,G),
\]
then $P$ is obviously representable as a closed sub-prescheme of finite presentation of $({}_nG)^r$, where
\[
{}_nG=\uHom_{S\text{-gr}}((\ZZ/n\ZZ)_S,G),
\]
which is also the kernel of the $n$th-power morphism in $G$, hence representable by a closed sub-prescheme of finite presentation of $G$. Moreover $P$ is smooth over $S$ by XI~2.1. Let $s\in S$; by what was seen above, there exists a finite separable extension $k''$ of $\kappa(s)$ such that there exists a maximal torus $T''$ in $G_{k''}$; in addition, since ${}_nT''$ is étale over $k''$ ($n$ being prime to the characteristic of $k''$), one can suppose (on replacing $k''$ by a finite separable extension) that ${}_nT''$ is isomorphic to $M_{k''}$. Let $S''\to S$ be an étale morphism such that there exists an $s''\in S''$ above $s\in S$ giving rise to the residue extension $\kappa(s'')\simeq k''$. One therefore has a section of $P_{S''}\otimes\kappa(s'')$ over $\Spec(\kappa(s''))$, so by Hensel, on replacing $(S'',s'')$ by $(S''',s''')$ étale over it, one can suppose that there exists a section of $P_{S''}$ over $S''$, \ie an element of $P(S'')$, extending the given section. In other words, one has a homomorphism $M_{S''}\to G_{S''}$ inducing an isomorphism $M_{\kappa(s'')}\simeq{}_nT''$. By IX~6.4 (which here reduces to a simple application of Nakayama's lemma), this homomorphism is a closed immersion over an open neighborhood of $s'$, which one can suppose equal to $S''$. Let $H''$ be its image, and consider its centralizer $C''$ in $G_{S''}$, which is a smooth closed sub-prescheme of $G$, by XI~6.2. Note moreover:
% typo?: unprimed C,T in C=Centr_G(nT), and s' instead of s'', kept as printed.
\begin{lemma}{7.3}\label{XII.7.3}
Under the preceding conditions for $\ell,n,r$, for every prescheme $S''$ over $S$ and every maximal torus $T''$ of $G_{S''}$, one has
\[
\Centr_{G_{S''}}(T'')=\Centr_{G_{S''}}({}_nT'').
\]
\end{lemma}
Indeed, by faithfully flat descent from $S'\times_S S''$, one is reduced to the case where $G_{S''}$ admits a maximal torus $T''_1$ for which the preceding relation is true. But since $T''$ is locally conjugate to $T''_1$ for the étale topology by b), it follows that the same relation is true for $T''$.

Apply the preceding result with $\Spec(\kappa(s''))$ instead of $S''$; one sees that the fiber $C''_{s''}$ is a Cartan subgroup of $G$. Note now that since $G$ is smooth over $S$, the union of the neutral components of the fibers $G_s$ is an open subset of $G$ (by a general result of EGA~IV on smooth morphisms\nde{Make this reference precise.}), obviously stable under the group law of $G$, hence is a subgroup of $G$ for the prescheme structure induced by $G$.

In addition, $G^0$ satisfies the preliminary hypotheses on $G$, and there is an obvious one-to-one correspondence between the maximal tori of $G$ and those of $G^0$. Thus to prove c) and d), one can suppose $G$ with connected fibers, which we shall do. Then the Cartan subgroups of $G$ have connected fibers (6.6 a)). This being said, I say that $C''{}^0$ is a Cartan subgroup of $G_{S''}$ over an open neighborhood of $s''$ in $S''$ (which will complete the proof of c)). This follows indeed from the
\begin{lemma}{7.4}\label{XII.7.4}
Under the conditions of 7.1, suppose $G$ with connected fibers; let $C$ be a closed group sub-prescheme of $G$, smooth over $S$, and $s$ an element of $S$ such that $C_s$ is a Cartan subgroup of $G_s$; then $C^0$ is a Cartan subgroup of $G$ over an open neighborhood of $s$.
\end{lemma}
One reduces easily to the case where $S$ is local and $s$ its closed point, and to proving that $C$ is then a Cartan subgroup of $G$, then by flat descent to the case where $G$ admits a maximal torus, say $T$. Then by XI~6.2, $\uTransp_G(T,C)$ is representable by a closed sub-prescheme of $G$ smooth over $S$. In addition, by the hypothesis on $C_s$, the fiber of the transporter at $s$ is nonempty (taking account of conjugacy theorem 6.6 a)). This reduces us by faithfully flat descent to the case where this transporter admits a section over $S$, hence to the case where $C$ contains a maximal torus $T$ of $G$. But then $T$ is central in $C^0$ by IX~5.6 a), hence $C^0\subset\Centr_G(T)$, and since this is an inclusion of group schemes smooth over $S$, having the same relative dimension (namely the dimension of their common fiber at $s$) and connected fibers, it is an equality, which completes the proof.
% typo?: proving C rather than C^0 at the beginning kept as printed.

d) We retain the preceding notation and hypotheses for $\ell,n,r$, and the connectedness of the fibers of $G$. Let $Q:\Sch^\circ_{/S}\to\Ens$ be the functor defined by
\[
\begin{split}
Q(S')=\text{set of subgroups of multiplicative type of }G_{S'}\\
\text{of type equal to }(\ZZ/n\ZZ)^r=M\quad\text{(IX 1.4).}
\end{split}
\]
Then
\[
T\mto{}_nT
\]
is a morphism
\[
\varphi:\mathscr T\to Q,
\]
which is a monomorphism by 7.3. I say that $Q$ is representable by a separated prescheme of finite presentation over $S$. Indeed, as was pointed out in the proof of XI~3.12 a), one has an isomorphism
\[
Q\simeq P'/\Gamma
\]
where $P'$ is the open and closed sub-prescheme of the prescheme $P=\uHom_{S\text{-gr}}(M_S,G)$ introduced in c) corresponding to monomorphisms $M_{S'}\to G_{S'}$ (\cf IX~6.8), and where $\Gamma=(\Aut_{\mathrm{gr}}(M))_S$. The hypothesis that every finite subset of a fiber $G_s$ is contained in an affine open subset of $G$, being stable under passage to a closed sub-prescheme and under cartesian products, is obviously inherited by $({}_nG)^r$, hence by $P$, hence by $P'$, so $Q$ is representable by a prescheme of finite presentation over $S$ (\cf V)\nde{Make this reference precise.}.

One sees in the same way that if $G$ is quasi-projective (resp. affine) over $S$, so is $Q$. In any case, $Q$ is separated over $S$. Now one has the
\begin{lemma}{7.5}\label{XII.7.5}
The preceding homomorphism $\varphi=\mathscr T\to Q$ is representable by an open immersion.
% typo?: equals sign rather than colon kept as printed.
\end{lemma}
In other words, one must prove that if $H$ is a subgroup of multiplicative type of $G$, of type $M$, then there exists an open subset $U$ of $S$ such that for every $S'$ over $S$, $H_{S'}$ is of the form ${}_nT_{S'}$, for a suitable maximal torus $T_{S'}$ of $G_{S'}$, if and only if $S'\to S$ factors through $U$. One can obviously suppose that the nilpotent rank of the fibers of $G$ is constant (for, thanks to b), it is locally constant), say $r'$. Let $C=\Centr_G(H)$, which is a closed group sub-prescheme of $G$ smooth over $S$ (XI~6.2). Then, on replacing $S$ by the open and closed subset of points at which $C$ is of relative dimension $r'$, one can suppose $C$ of relative dimension $r'$ everywhere. One then sees at once that $H$ is of the form ${}_nT$, for a maximal torus $T$ of $G$, if and only if $C^0$ admits a central maximal torus $T$ of relative dimension $r$ everywhere and $H={}_nT$, which gives another expression for the subfunctor $U$ of $S$ that we want to represent (replacing $S$ in the preceding criterion by an $S'$ over $S$). Moreover, by flat descent one can suppose that $G$ admits a maximal torus $T_1$. Let $R$ be the subfunctor $\uTransp_G(T_1,\Centr_{C^0}(C^0))$ of $\uTransp_G(T_1,C)$; the latter is representable by a smooth prescheme over $S$ by XI~6.2, and the former is representable by an induced open sub-prescheme, as follows at once from IX~5.6 a); in particular it is smooth over $S$. Consequently its structural morphism to $S$ is open, hence its image is open, and on replacing $S$ by the said image (endowed with the induced structure), one can suppose the structural morphism surjective. Then, by 1.13 applied to $C^0$, $C^0$ admits a central maximal torus $T$, since it admits one locally for fpqc, which will obviously be of relative dimension equal to that of $T_1$, \ie $r$. Thus the condition to express for $H$ is the equality $H={}_nT$, which by IX~2.10 amounts again to taking a suitable open (and closed) subset of $S$.

Lemma 7.5 therefore implies that $\mathscr T$ is representable by a separated prescheme locally of finite presentation over $S$, and even of finite presentation over $S$, as one sees by taking up again the proof of 7.5 to ensure that $\varphi$ is in fact a quasi-compact open immersion, or by reducing by faithfully flat quasi-compact descent to the case where $G$ admits a maximal torus, and where $\mathscr T$ is therefore isomorphic to $G/\Norm_G(T)$. This last expression, or alternatively XI~2.1, shows in addition that $\mathscr T$ is smooth over $S$. Finally, if $G$ is quasi-projective over $S$, so is $Q$, hence also $\mathscr T$. If $G$ is affine over $S$, the assertion that $\mathscr T$ is then affine over $S$ is included as a reminder, having been established in 5.4 (N.B. I do not know whether, without an affine hypothesis for $G/S$, it is possible to choose $n$ so that in 7.5 the open immersion is also a closed immersion). For the assertion that $\mathscr T$ is affine over $S$ if $S$ is artinian, one is reduced to the case where $S$ is the spectrum of a field (EGA~I~6.1.7), which one can suppose algebraically closed. Then thanks to f), which will be proved below, it suffices to prove the same assertion for $G/\Centr(G)$; but the latter is affine by 6.1, so one is under the preceding conditions.

This completes the proof of d).

e) By IX~6.8, one knows that there exists a subtorus $T'$ of $G'$ such that $u$ induces a faithfully flat morphism $T\to T'$ (which characterizes $T'$ as the subsheaf $u(T)$ of $G'$). Let $C$ be the centralizer of $T$, $C'$ that of $T'$; let us prove that the morphism $C\to C'$ is flat, and faithfully flat if $G'$ has connected fibers. Since $C,C'$ are flat of finite presentation over $S$, one is reduced to the case of a base field (SGA~1 I~5.9), which one can obviously suppose algebraically closed. In addition one can suppose $G,G'$ connected (on replacing them by $G^0$ and $G'{}^0$, which does not change $C^0$ and $C'{}^0$), and it then suffices to apply 6.6 d), taking account of a).

f) Taking account of a) and e), one can restrict oneself to proving the assertion concerning Cartan subgroups. Now since a Cartan subgroup of $G$ is the centralizer of a maximal torus, it contains the center of $G$ and a fortiori $\Ker u$, so it is of the form $u^{-1}(C')$, where $C'=u(C)$ is the Cartan subgroup of $G'$ considered in e). Thus the map $C\mto u(C)$ is injective; to show that it is bijective, it suffices to see that for every Cartan subgroup $C'$ of $G'$, $u^{-1}(C')$ is a Cartan subgroup of $G$. Since the question is local for fpqc, one can suppose that $G$ admits a Cartan subgroup $C_1$, hence $u(C_1)=C'_1$ is a Cartan subgroup of $G'$, hence locally conjugate to $C'$ for fpqc by b), and since $u^{-1}(C'_1)=C_1$ is a Cartan subgroup of $G$, it follows that $u^{-1}(C')$ is a Cartan subgroup of $G$. \hfill Q.E.D.

One can also, to prove that $u^{-1}(C')$ is a Cartan subgroup, note that it is flat over $S$, for $u$ is (SGA~1 I~5.9), which reduces us by definition to the case of a base field, and one can apply 6.6 e).
\end{proof}
\begin{corollary}{7.6}\label{XII.7.6}
Under the conditions of 7.1 e), $G'$ admits locally for the étale topology a maximal torus, hence satisfies the preliminary conditions for $G$. If in addition $\Ker u$ is central, then the functors $\mathscr T_G,\mathscr C_G$ of maximal tori of $G$ and Cartan subgroups of $G$ are isomorphic to the analogous functors $\mathscr T_{G'},\mathscr C_{G'}$ for $G'$ (hence, in the case where they are representable, they are represented by isomorphic $S$-preschemes).
\end{corollary}
\begin{remark}{7.7}\label{XII.7.7}
a) Unlike what happens in the case where $G$ is affine over $S$ (it suffices, in fact, that $G$ have affine fibers, as one will see in XVI), it is not true that the fact that $G$ has locally constant reductive rank implies that $G$ admits locally for fpqc a maximal torus. Let, for example, $S$ be the spectrum of a discrete valuation ring, $G_1$ and $G_2$ smooth separated group schemes of finite type over $S$, such that the generic fiber of $G_1$ is an elliptic curve, the special fiber a group $\Gm$, and the generic fiber of $G_2$ a group $\Gm$, the special fiber a group $\Ga$, and take $G=G_1\times_S G_2$. Then the two fibers of $G$ have reductive rank~1, but one sees at once that $G$ admits no maximal torus locally for fpqc. On the other hand, it is very plausible that the following condition (for a smooth separated group of finite type over a prescheme $S$) is sufficient for the existence of a maximal torus locally for the étale topology: the reductive rank and the abelian rank of the fibers of $G$ are locally constant functions.

b) In the proof of 7.1 (in particular a)), we invoked XI~6.2 in cases where $S$ is not supposed locally noetherian. However, in the cases considered of application of XI~6.2, reduction to the case of affine noetherian $S$ is immediate.
\end{remark}
Here is a variant of 7.1 b):
\begin{proposition}{7.8}\label{XII.7.8}
Let $G$ be a smooth $S$-group prescheme of finite presentation with connected fibers, $H$ a group prescheme smooth over $S$, $i:H\to G$ a monomorphism of $S$-groups (making $H$ a sub-$S$-group of $G$). Then for every Cartan subgroup{\renewcommand{\thefootnote}{*}\footnote{The proof shows that it suffices to suppose that $C$ is a smooth subgroup of $G$ each of whose geometric fibers is the connected centralizer of a subgroup of a maximal torus of $G_{\overline S}$.}\addtocounter{footnote}{-1}} $C$ of $G$, $\uTransp_G(C,H)$ is representable by a closed sub-prescheme of $G$ smooth over $S$.
% typo?: capital S in G_{bar S} in the starred note kept as printed.
\end{proposition}
The fact that the transporter is representable by a closed sub-prescheme of $G$ of finite presentation is contained in XI~6.11, taking into account that $C$ has connected fibers since $G$ does (6.6 b)). To show that the transporter is smooth over $S$, one is reduced by the standard procedure to the case where $S$ is affine noetherian, then to the case where $S$ is artinian local, and by descent to the case where the residue field of $S$ is algebraically closed. But then $C$ admits a maximal torus $T$, which is a maximal torus of $G$. One can suppose that the reductive rank and the nilpotent rank of the fiber $H_0$ are equal to those of $G_0$ (otherwise the transporter would be empty), but then one sees at once (using the connectedness of $C$ and the fact that the centralizer in $H$ of a maximal torus of $H$ is smooth) that one has
\[
\uTransp_G(C,H)=\uTransp_G(T,H)
\]
and since one knows that the second member is smooth (XI~2.5), so is the first. The preceding argument proves more generally part b) of the
\begin{proposition}{7.9}\label{XII.7.9}
Let $G$ and $i:H\to G$ be as in 7.8; suppose in addition that for every $s\in S$, $H_s$ is connected, and has the same reductive rank and the same nilpotent rank as $G_s$ (\ie $H_s$ contains a Cartan subgroup of $G_s$). Then one has the following:
\begin{enumerate}[label=\alph*)]
\item $\uNorm_G(H)$ is representable by a closed sub-prescheme $\Norm_G(H)$ of $G$ of finite presentation over $S$, and the canonical monomorphism $H\to\Norm_G(H)$ is an open immersion; consequently $i$ is an immersion, and one has
\[
H=\Norm_G(H)^0.
\]
\item For every Cartan subgroup $C$ of $G$, $\uTransp_G(C,H)$ is a closed sub-prescheme of $G$, smooth over $S$, with surjective structural morphism. If $C$ is the centralizer of a maximal torus $T$ of $G$, one has in addition
\[
\uTransp_G(T,H)=\uTransp_G(C,H).
\]
\item Let $C$ be a group sub-prescheme of $H$. For it to be a Cartan subgroup of $H$, it is necessary and sufficient that it be a Cartan subgroup of $G$.
\item Suppose that $G$ admits locally for the étale topology, or for the fpqc topology, a Cartan subgroup (resp. a maximal torus); then the same holds for $H$.
\end{enumerate}
\end{proposition}
\begin{proof}
a) The representability of $\uNorm_G(H)$ by a closed sub-prescheme $\Norm_G(H)$ of $G$ of finite presentation over $S$ is contained in XI~6.11. Since $H$ is smooth, hence flat locally of finite presentation over $S$, to verify that $H\to\Norm_G(H)$ is an open immersion, one is reduced to verifying it on the fibers (VI$_{\mathrm B}$.2.6), which reduces us to the case where $S$ is the spectrum of an algebraically closed field $k$. One is then reduced (Exp.~VI\nde{Reference not located in VI$_{\mathrm B}$; see M.~Demazure and P.~Gabriel, \textit{Groupes algébriques}, I, Masson (1970), corollary II.5.6.}) to verifying that the corresponding homomorphism on the Lie algebras is an isomorphism, or, what amounts to the same thing, that $(\mathfrak g/\mathfrak h)^H=0$, where $\mathfrak g$ and $\mathfrak h$ are the Lie algebras of $G$ and $H$, and the exponent $H$ denotes the invariants under $H$ (\cf II~5.2.3 (i)). Now by hypothesis $H$ contains a Cartan subgroup $C$ of $G$, the centralizer of the maximal torus $T$ of $G$, and it therefore suffices to prove that one has
\[
(\mathfrak g/\mathfrak h)^T=0,
\]
which follows, taking account of the complete reducibility of representations of $T$ (I~4.7.3), from the analogous relation $(\mathfrak g/\mathfrak c)^T$, where $\mathfrak c=\Lie(C)$. As for the latter, equivalent to
\[
\mathfrak g^T=\mathfrak c,
\]
it means that the centralizer $C$ and the normalizer $N$ of $T$ have the same Lie algebra, which follows from the fact that $C$ is an open subgroup of $N$ (XI~5.9). This completes the proof of a).
% typo?: the relation (g/c)^T is printed without =0, retained.

b) As was pointed out, the proof was given in 7.8.

c) Taking account of the fact that $H$ is a sub-prescheme of $G$, the assertion reduces trivially to the case where $S$ is the spectrum of an algebraically closed field, in which case it follows at once from the hypothesis made on $H$.

d) One uses b), c) and ``Hensel's lemma'' XI~1.10.
\end{proof}
\begin{corollary}{7.10}\label{XII.7.10}
Let $G$ and $H$ be as in 7.9, and suppose that for every algebraically closed field $k$ over $S$, every element of $G_k(k)$ normalizing $H_k$ is in $H_k(k)$. Then $H$ is a closed sub-prescheme of $G$, and is its own normalizer.
\end{corollary}
This follows trivially from 7.9 a). We shall apply 7.10 in particular to the Borel subgroups (more generally, to the parabolic subgroups) of $G$.
\begin{corollary}{7.11}\label{XII.7.11}
Let $G,H$ be as in 7.9. Then $H$ contains every group sub-prescheme $Z$ of $G$, central in $G$ and flat and of finite presentation over $S$.
\end{corollary}
One can as usual reduce to the case of affine noetherian $S$, then to the case of artinian $S$, which implies that one is under the conditions of 7.1. By 7.9 c), one can suppose that $H$ contains a Cartan subgroup $C$ of $G$, and one is reduced to proving that $C\supset Z$. Now since $S$ is artinian, $G'=G/Z$ is representable by a group prescheme of finite type over $S$, the canonical morphism $u:G\to G'$ being faithfully flat and its kernel being $Z$ (VI$_{\mathrm A}$.3.2). Obviously $G'$ is smooth over $S$, and one can apply 7.1 f), which implies that $C$ is of the form $u^{-1}(C')$, hence contains $Z$.
